\documentclass[11pt,a4paper]{article}

\def\courseTitle{Industrial Security (Master)}
\def\labNumber{06}
\def\labTitle{Threat Intelligence for OT -- Solutions}
\def\labSection{06}
\def\labTime{3--4 hours}

% =====================================================================
% Shared preamble for lab sheets.
%
% Caller must \documentclass{article} first, then define:
%   \def\courseTitle{...}
%   \def\labNumber{NN}
%   \def\labTitle{...}
%   \def\labTime{e.g. "30--45 min"}
%   \def\labSection{e.g. "02"}
% Then % =====================================================================
% Shared preamble for lab sheets.
%
% Caller must \documentclass{article} first, then define:
%   \def\courseTitle{...}
%   \def\labNumber{NN}
%   \def\labTitle{...}
%   \def\labTime{e.g. "30--45 min"}
%   \def\labSection{e.g. "02"}
% Then % =====================================================================
% Shared preamble for lab sheets.
%
% Caller must \documentclass{article} first, then define:
%   \def\courseTitle{...}
%   \def\labNumber{NN}
%   \def\labTitle{...}
%   \def\labTime{e.g. "30--45 min"}
%   \def\labSection{e.g. "02"}
% Then \input{../../template/lab-preamble.tex}
% =====================================================================

\usepackage[utf8]{inputenc}
\usepackage[T1]{fontenc}
\usepackage[english]{babel}
\usepackage[a4paper, margin=2cm]{geometry}
\usepackage{graphicx}
\usepackage{listings}
\usepackage{xcolor}
\usepackage{colortbl}
\usepackage{tikz}
\usepackage{amsmath}
\usepackage{amssymb}
\usepackage{booktabs}
\usepackage{enumitem}
\usepackage{titlesec}
\usepackage{fancyhdr}
\usepackage{hyperref}
\usepackage{tcolorbox}
\tcbuselibrary{skins, breakable}
\usepackage{fontawesome5}

\input{../../../template/tikz-styles.tex}

\providecommand{\courseAuthor}{Industrial Security Lecture Series}
\providecommand{\courseInstitute}{Department of Computer Science}
\providecommand{\companyLogo}{logo}

\definecolor{slideAccent}{HTML}{00205B}
\definecolor{slideSecondary}{HTML}{00A0DC}

% --- Corporate branding overrides ------------------------------------
\IfFileExists{../../../branding.tex}{\input{../../../branding.tex}}{}

\colorlet{labAccent}{slideAccent}
\colorlet{labSec}{slideSecondary}
\definecolor{labCheck}{HTML}{2E7D32}
\definecolor{labWarn}{HTML}{B23A48}

\lstset{
  backgroundcolor=\color{black!4},
  basicstyle=\ttfamily\footnotesize,
  breaklines=true,
  frame=leftline,
  framerule=2pt,
  rulecolor=\color{labAccent},
  numbers=left,
  numberstyle=\tiny\color{black!50},
  showstringspaces=false,
  tabsize=2
}

\setlength{\tabcolsep}{4pt}

% Let TeX add a little extra inter-word stretch as a last resort before
% it reports an Overfull \hbox. Only kicks in on lines that would
% otherwise overflow (long \texttt paths, URLs, CIDRs); never loosens a
% line that already fits. Keeps prose/itemize within the margin.
\setlength{\emergencystretch}{3em}

\pagestyle{fancy}
\fancyhf{}
\renewcommand{\headrulewidth}{0.4pt}
\lhead{\small \courseTitle{} -- Lab \labNumber}
\rhead{\small Maps to Section \labSection}
\cfoot{\thepage}

\newtcolorbox{stepbox}[1]{
  enhanced, breakable, arc=2pt, boxrule=0pt,
  colback=labAccent!7, colframe=labAccent, leftrule=3pt,
  fonttitle=\bfseries\small, coltitle=white,
  colbacktitle=labAccent,
  title={\faTools\ Step #1}
}

\newtcolorbox{warnbox}{
  enhanced, breakable, arc=2pt, boxrule=0pt,
  colback=labWarn!7, colframe=labWarn, leftrule=4pt,
  fonttitle=\bfseries\small, coltitle=white, colbacktitle=labWarn,
  title={\faExclamationTriangle\ Caution}
}

\newtcolorbox{notebox}{
  enhanced, breakable, arc=2pt, boxrule=0pt,
  colback=labAccent!4, colframe=labAccent, leftrule=2pt,
  fonttitle=\bfseries\small, coltitle=white, colbacktitle=labAccent,
  title={\faInfoCircle\ Note}
}

\newtcolorbox{goalbox}{
  enhanced, breakable, arc=2pt, boxrule=0pt,
  colback=labCheck!7, colframe=labCheck, leftrule=3pt,
  fonttitle=\bfseries\small, coltitle=white, colbacktitle=labCheck,
  title={\faBullseye\ Goal}
}

\newtcolorbox{deliverbox}{
  enhanced, breakable, arc=2pt, boxrule=0pt,
  colback=labSec!10, colframe=labSec, leftrule=3pt,
  fonttitle=\bfseries\small, coltitle=white, colbacktitle=labSec,
  title={\faClipboardList\ Deliverable}
}

% --- Solution-sheet boxes (used only by lab-solutions.tex) ----------
\newtcolorbox{solutionbox}[1]{
  enhanced, breakable, arc=2pt, boxrule=0pt,
  colback=labCheck!7, colframe=labCheck, leftrule=3pt,
  fonttitle=\bfseries\small, coltitle=white, colbacktitle=labCheck,
  title={\faCheckCircle\ #1}
}

\newtcolorbox{rubricbox}{
  enhanced, breakable, arc=2pt, boxrule=0pt,
  colback=labSec!7, colframe=labSec, leftrule=3pt,
  fonttitle=\bfseries\small, coltitle=white, colbacktitle=labSec,
  title={\faBalanceScale\ Grading rubric}
}

\titleformat{\section}{\large\bfseries\color{labAccent}}{\thesection}{0.5em}{}

\newcommand{\labheader}{%
  \noindent
  \begin{tikzpicture}
    \fill[labAccent] (0,0) rectangle (\textwidth, -1.6cm);
    \node[anchor=north west, text=white, font=\bfseries\large] at (0.6cm, -0.4cm) {Lab \labNumber{} -- \labTitle};
    \node[anchor=south west, text=white, font=\small] at (0.6cm, -1.3cm) {\courseTitle{} \quad $\bullet$ \quad Maps to Section \labSection{} \quad $\bullet$ \quad \labTime};
    \node[anchor=north east, fill=labSec, text=white, font=\bfseries, inner sep=4pt, rounded corners=2pt]
      at ([xshift=-0.4cm, yshift=-0.4cm]\textwidth,0) {LAB \labNumber};
  \end{tikzpicture}
  \vspace{4mm}
}

% =====================================================================

\usepackage[utf8]{inputenc}
\usepackage[T1]{fontenc}
\usepackage[english]{babel}
\usepackage[a4paper, margin=2cm]{geometry}
\usepackage{graphicx}
\usepackage{listings}
\usepackage{xcolor}
\usepackage{colortbl}
\usepackage{tikz}
\usepackage{amsmath}
\usepackage{amssymb}
\usepackage{booktabs}
\usepackage{enumitem}
\usepackage{titlesec}
\usepackage{fancyhdr}
\usepackage{hyperref}
\usepackage{tcolorbox}
\tcbuselibrary{skins, breakable}
\usepackage{fontawesome5}

% =====================================================================
% Shared TikZ styles for the Industrial Security lecture series.
% Loaded by every slide deck and exercise sheet that uses TikZ.
% =====================================================================

\usetikzlibrary{
  arrows.meta,
  shapes,
  shapes.geometric,
  shapes.symbols,
  shapes.multipart,
  positioning,
  fit,
  calc,
  backgrounds,
  decorations.pathmorphing,
  decorations.pathreplacing,
  decorations.text,
  patterns,
  matrix,
  shadows
}

% --- colour palette --------------------------------------------------
\definecolor{isOT}{HTML}{1F4E79}     % OT (deep blue)
\definecolor{isIT}{HTML}{6F2DA8}     % IT (purple)
\definecolor{isSafety}{HTML}{C00000} % safety red
\definecolor{isNet}{HTML}{2E7D32}    % network green
\definecolor{isAttack}{HTML}{B23A48} % attack red
\definecolor{isAccent}{HTML}{E08E0B} % amber
\definecolor{isMute}{HTML}{6B6B6B}   % grey
\definecolor{isLight}{HTML}{EDF2F8}  % light background
\definecolor{isPLC}{HTML}{2C5282}
\definecolor{isHMI}{HTML}{2F855A}
\definecolor{isSCADA}{HTML}{6B46C1}

% --- generic node styles --------------------------------------------
\tikzset{
  isnode/.style={
    draw, rounded corners=2pt, minimum height=8mm, minimum width=18mm,
    font=\small, align=center, thick
  },
  isbox/.style={isnode, fill=isLight},
  plc/.style={isnode, fill=isPLC!15, draw=isPLC, thick},
  hmi/.style={isnode, fill=isHMI!15, draw=isHMI, thick},
  scada/.style={isnode, fill=isSCADA!15, draw=isSCADA, thick},
  sensor/.style={isnode, fill=isNet!10, draw=isNet, minimum height=6mm, minimum width=14mm, font=\scriptsize},
  actuator/.style={isnode, fill=isAccent!15, draw=isAccent, minimum height=6mm, minimum width=14mm, font=\scriptsize},
  firewall/.style={isnode, fill=isAttack!15, draw=isAttack, thick},
  server/.style={isnode, fill=isMute!15, draw=isMute!80, thick},
  switch/.style={isnode, fill=isLight, draw=black!60, minimum height=6mm, minimum width=14mm, font=\scriptsize},
  workstation/.style={isnode, fill=white, draw=isOT, thick},
  attacker/.style={isnode, fill=isAttack!20, draw=isAttack, thick, font=\small\bfseries},
  inetnode/.style={draw, ellipse, minimum width=24mm, minimum height=10mm, font=\scriptsize, fill=isLight, thick},
  process/.style={isnode, fill=isOT!15, draw=isOT, thick, minimum width=24mm},
  zone/.style={
    draw, rounded corners=4pt, thick, dashed, inner sep=4mm
  },
  conduit/.style={-{Stealth[length=2.5mm]}, thick, isNet!80!black},
  attack/.style={-{Stealth[length=2.5mm]}, thick, isAttack, dashed},
  flow/.style={-{Stealth[length=2.5mm]}, thick, isOT!70!black},
  netline/.style={thick, black!60},
  axisline/.style={-{Stealth[length=2mm]}, thick, black!70},
  tag/.style={font=\scriptsize\itshape, fill=white, inner sep=1pt},
  level/.style={
    draw, fill=isLight, minimum width=0.85\linewidth, minimum height=8mm,
    font=\small, anchor=west
  },
  brace/.style={decorate, decoration={brace, amplitude=4pt, raise=2pt}},
  legendbox/.style={draw, rounded corners=2pt, fill=isLight, inner sep=2mm, font=\scriptsize, align=left}
}

% --- handy macros ----------------------------------------------------
% \plcicon, \hmiicon etc as simple labelled boxes; useful inline.
\newcommand{\plcicon}[1][PLC]{\tikz[baseline=-0.6ex]{\node[plc, minimum width=10mm, minimum height=5mm, font=\scriptsize, inner sep=1pt]{#1};}}
\newcommand{\hmiicon}[1][HMI]{\tikz[baseline=-0.6ex]{\node[hmi, minimum width=10mm, minimum height=5mm, font=\scriptsize, inner sep=1pt]{#1};}}
\newcommand{\fwicon}[1][FW]{\tikz[baseline=-0.6ex]{\node[firewall, minimum width=10mm, minimum height=5mm, font=\scriptsize, inner sep=1pt]{#1};}}


\providecommand{\courseAuthor}{Industrial Security Lecture Series}
\providecommand{\courseInstitute}{Department of Computer Science}
\providecommand{\companyLogo}{logo}

\definecolor{slideAccent}{HTML}{00205B}
\definecolor{slideSecondary}{HTML}{00A0DC}

% --- Corporate branding overrides ------------------------------------
\IfFileExists{../../../branding.tex}{% =====================================================================
% Corporate / institutional branding -- single file to customise the
% look of the entire course (slides, notes, exercises, labs).
%
% HOW TO REBRAND IN UNDER A MINUTE:
%   1. Replace `branding/logo.pdf` with your own logo
%      (PDF preferred; PNG and JPG also work -- name it `logo.<ext>`).
%   2. (Optional) change the two colours below to your brand palette.
%   3. (Optional) override the author/institute lines below.
%   4. Re-run `make` at the repo root. That's it.
%
% Everything in this file has sensible defaults; you can delete a line
% to fall back to the built-in defaults, or comment it out.
%
% This file is loaded automatically by the slides, exercise, and lab
% preambles -- you never need to \input it manually.
% =====================================================================

% --- Logo --------------------------------------------------------------
% Path is resolved from each leaf-build directory; the default points at
% the shared `branding/` folder at the repo root. If you put your logo
% somewhere else, set the full or relative path here (without extension):
\renewcommand{\companyLogo}{../../../branding/logo}

% --- Author / institute (appears on the title page footer) ------------
\renewcommand{\courseAuthor}{}
\renewcommand{\courseInstitute}{}

% --- Brand colours ----------------------------------------------------
% slideAccent      = primary brand colour (title bands, accent rule)
% slideSecondary   = secondary brand colour (section badge, highlight)
% Both use HTML hex codes. Keep enough contrast against white text.
\definecolor{slideAccent}{HTML}{00205B}      % deep navy
\definecolor{slideSecondary}{HTML}{00A0DC}   % light blue accent
}{}

\colorlet{labAccent}{slideAccent}
\colorlet{labSec}{slideSecondary}
\definecolor{labCheck}{HTML}{2E7D32}
\definecolor{labWarn}{HTML}{B23A48}

\lstset{
  backgroundcolor=\color{black!4},
  basicstyle=\ttfamily\footnotesize,
  breaklines=true,
  frame=leftline,
  framerule=2pt,
  rulecolor=\color{labAccent},
  numbers=left,
  numberstyle=\tiny\color{black!50},
  showstringspaces=false,
  tabsize=2
}

\setlength{\tabcolsep}{4pt}

% Let TeX add a little extra inter-word stretch as a last resort before
% it reports an Overfull \hbox. Only kicks in on lines that would
% otherwise overflow (long \texttt paths, URLs, CIDRs); never loosens a
% line that already fits. Keeps prose/itemize within the margin.
\setlength{\emergencystretch}{3em}

\pagestyle{fancy}
\fancyhf{}
\renewcommand{\headrulewidth}{0.4pt}
\lhead{\small \courseTitle{} -- Lab \labNumber}
\rhead{\small Maps to Section \labSection}
\cfoot{\thepage}

\newtcolorbox{stepbox}[1]{
  enhanced, breakable, arc=2pt, boxrule=0pt,
  colback=labAccent!7, colframe=labAccent, leftrule=3pt,
  fonttitle=\bfseries\small, coltitle=white,
  colbacktitle=labAccent,
  title={\faTools\ Step #1}
}

\newtcolorbox{warnbox}{
  enhanced, breakable, arc=2pt, boxrule=0pt,
  colback=labWarn!7, colframe=labWarn, leftrule=4pt,
  fonttitle=\bfseries\small, coltitle=white, colbacktitle=labWarn,
  title={\faExclamationTriangle\ Caution}
}

\newtcolorbox{notebox}{
  enhanced, breakable, arc=2pt, boxrule=0pt,
  colback=labAccent!4, colframe=labAccent, leftrule=2pt,
  fonttitle=\bfseries\small, coltitle=white, colbacktitle=labAccent,
  title={\faInfoCircle\ Note}
}

\newtcolorbox{goalbox}{
  enhanced, breakable, arc=2pt, boxrule=0pt,
  colback=labCheck!7, colframe=labCheck, leftrule=3pt,
  fonttitle=\bfseries\small, coltitle=white, colbacktitle=labCheck,
  title={\faBullseye\ Goal}
}

\newtcolorbox{deliverbox}{
  enhanced, breakable, arc=2pt, boxrule=0pt,
  colback=labSec!10, colframe=labSec, leftrule=3pt,
  fonttitle=\bfseries\small, coltitle=white, colbacktitle=labSec,
  title={\faClipboardList\ Deliverable}
}

% --- Solution-sheet boxes (used only by lab-solutions.tex) ----------
\newtcolorbox{solutionbox}[1]{
  enhanced, breakable, arc=2pt, boxrule=0pt,
  colback=labCheck!7, colframe=labCheck, leftrule=3pt,
  fonttitle=\bfseries\small, coltitle=white, colbacktitle=labCheck,
  title={\faCheckCircle\ #1}
}

\newtcolorbox{rubricbox}{
  enhanced, breakable, arc=2pt, boxrule=0pt,
  colback=labSec!7, colframe=labSec, leftrule=3pt,
  fonttitle=\bfseries\small, coltitle=white, colbacktitle=labSec,
  title={\faBalanceScale\ Grading rubric}
}

\titleformat{\section}{\large\bfseries\color{labAccent}}{\thesection}{0.5em}{}

\newcommand{\labheader}{%
  \noindent
  \begin{tikzpicture}
    \fill[labAccent] (0,0) rectangle (\textwidth, -1.6cm);
    \node[anchor=north west, text=white, font=\bfseries\large] at (0.6cm, -0.4cm) {Lab \labNumber{} -- \labTitle};
    \node[anchor=south west, text=white, font=\small] at (0.6cm, -1.3cm) {\courseTitle{} \quad $\bullet$ \quad Maps to Section \labSection{} \quad $\bullet$ \quad \labTime};
    \node[anchor=north east, fill=labSec, text=white, font=\bfseries, inner sep=4pt, rounded corners=2pt]
      at ([xshift=-0.4cm, yshift=-0.4cm]\textwidth,0) {LAB \labNumber};
  \end{tikzpicture}
  \vspace{4mm}
}

% =====================================================================

\usepackage[utf8]{inputenc}
\usepackage[T1]{fontenc}
\usepackage[english]{babel}
\usepackage[a4paper, margin=2cm]{geometry}
\usepackage{graphicx}
\usepackage{listings}
\usepackage{xcolor}
\usepackage{colortbl}
\usepackage{tikz}
\usepackage{amsmath}
\usepackage{amssymb}
\usepackage{booktabs}
\usepackage{enumitem}
\usepackage{titlesec}
\usepackage{fancyhdr}
\usepackage{hyperref}
\usepackage{tcolorbox}
\tcbuselibrary{skins, breakable}
\usepackage{fontawesome5}

% =====================================================================
% Shared TikZ styles for the Industrial Security lecture series.
% Loaded by every slide deck and exercise sheet that uses TikZ.
% =====================================================================

\usetikzlibrary{
  arrows.meta,
  shapes,
  shapes.geometric,
  shapes.symbols,
  shapes.multipart,
  positioning,
  fit,
  calc,
  backgrounds,
  decorations.pathmorphing,
  decorations.pathreplacing,
  decorations.text,
  patterns,
  matrix,
  shadows
}

% --- colour palette --------------------------------------------------
\definecolor{isOT}{HTML}{1F4E79}     % OT (deep blue)
\definecolor{isIT}{HTML}{6F2DA8}     % IT (purple)
\definecolor{isSafety}{HTML}{C00000} % safety red
\definecolor{isNet}{HTML}{2E7D32}    % network green
\definecolor{isAttack}{HTML}{B23A48} % attack red
\definecolor{isAccent}{HTML}{E08E0B} % amber
\definecolor{isMute}{HTML}{6B6B6B}   % grey
\definecolor{isLight}{HTML}{EDF2F8}  % light background
\definecolor{isPLC}{HTML}{2C5282}
\definecolor{isHMI}{HTML}{2F855A}
\definecolor{isSCADA}{HTML}{6B46C1}

% --- generic node styles --------------------------------------------
\tikzset{
  isnode/.style={
    draw, rounded corners=2pt, minimum height=8mm, minimum width=18mm,
    font=\small, align=center, thick
  },
  isbox/.style={isnode, fill=isLight},
  plc/.style={isnode, fill=isPLC!15, draw=isPLC, thick},
  hmi/.style={isnode, fill=isHMI!15, draw=isHMI, thick},
  scada/.style={isnode, fill=isSCADA!15, draw=isSCADA, thick},
  sensor/.style={isnode, fill=isNet!10, draw=isNet, minimum height=6mm, minimum width=14mm, font=\scriptsize},
  actuator/.style={isnode, fill=isAccent!15, draw=isAccent, minimum height=6mm, minimum width=14mm, font=\scriptsize},
  firewall/.style={isnode, fill=isAttack!15, draw=isAttack, thick},
  server/.style={isnode, fill=isMute!15, draw=isMute!80, thick},
  switch/.style={isnode, fill=isLight, draw=black!60, minimum height=6mm, minimum width=14mm, font=\scriptsize},
  workstation/.style={isnode, fill=white, draw=isOT, thick},
  attacker/.style={isnode, fill=isAttack!20, draw=isAttack, thick, font=\small\bfseries},
  inetnode/.style={draw, ellipse, minimum width=24mm, minimum height=10mm, font=\scriptsize, fill=isLight, thick},
  process/.style={isnode, fill=isOT!15, draw=isOT, thick, minimum width=24mm},
  zone/.style={
    draw, rounded corners=4pt, thick, dashed, inner sep=4mm
  },
  conduit/.style={-{Stealth[length=2.5mm]}, thick, isNet!80!black},
  attack/.style={-{Stealth[length=2.5mm]}, thick, isAttack, dashed},
  flow/.style={-{Stealth[length=2.5mm]}, thick, isOT!70!black},
  netline/.style={thick, black!60},
  axisline/.style={-{Stealth[length=2mm]}, thick, black!70},
  tag/.style={font=\scriptsize\itshape, fill=white, inner sep=1pt},
  level/.style={
    draw, fill=isLight, minimum width=0.85\linewidth, minimum height=8mm,
    font=\small, anchor=west
  },
  brace/.style={decorate, decoration={brace, amplitude=4pt, raise=2pt}},
  legendbox/.style={draw, rounded corners=2pt, fill=isLight, inner sep=2mm, font=\scriptsize, align=left}
}

% --- handy macros ----------------------------------------------------
% \plcicon, \hmiicon etc as simple labelled boxes; useful inline.
\newcommand{\plcicon}[1][PLC]{\tikz[baseline=-0.6ex]{\node[plc, minimum width=10mm, minimum height=5mm, font=\scriptsize, inner sep=1pt]{#1};}}
\newcommand{\hmiicon}[1][HMI]{\tikz[baseline=-0.6ex]{\node[hmi, minimum width=10mm, minimum height=5mm, font=\scriptsize, inner sep=1pt]{#1};}}
\newcommand{\fwicon}[1][FW]{\tikz[baseline=-0.6ex]{\node[firewall, minimum width=10mm, minimum height=5mm, font=\scriptsize, inner sep=1pt]{#1};}}


\providecommand{\courseAuthor}{Industrial Security Lecture Series}
\providecommand{\courseInstitute}{Department of Computer Science}
\providecommand{\companyLogo}{logo}

\definecolor{slideAccent}{HTML}{00205B}
\definecolor{slideSecondary}{HTML}{00A0DC}

% --- Corporate branding overrides ------------------------------------
\IfFileExists{../../../branding.tex}{% =====================================================================
% Corporate / institutional branding -- single file to customise the
% look of the entire course (slides, notes, exercises, labs).
%
% HOW TO REBRAND IN UNDER A MINUTE:
%   1. Replace `branding/logo.pdf` with your own logo
%      (PDF preferred; PNG and JPG also work -- name it `logo.<ext>`).
%   2. (Optional) change the two colours below to your brand palette.
%   3. (Optional) override the author/institute lines below.
%   4. Re-run `make` at the repo root. That's it.
%
% Everything in this file has sensible defaults; you can delete a line
% to fall back to the built-in defaults, or comment it out.
%
% This file is loaded automatically by the slides, exercise, and lab
% preambles -- you never need to \input it manually.
% =====================================================================

% --- Logo --------------------------------------------------------------
% Path is resolved from each leaf-build directory; the default points at
% the shared `branding/` folder at the repo root. If you put your logo
% somewhere else, set the full or relative path here (without extension):
\renewcommand{\companyLogo}{../../../branding/logo}

% --- Author / institute (appears on the title page footer) ------------
\renewcommand{\courseAuthor}{}
\renewcommand{\courseInstitute}{}

% --- Brand colours ----------------------------------------------------
% slideAccent      = primary brand colour (title bands, accent rule)
% slideSecondary   = secondary brand colour (section badge, highlight)
% Both use HTML hex codes. Keep enough contrast against white text.
\definecolor{slideAccent}{HTML}{00205B}      % deep navy
\definecolor{slideSecondary}{HTML}{00A0DC}   % light blue accent
}{}

\colorlet{labAccent}{slideAccent}
\colorlet{labSec}{slideSecondary}
\definecolor{labCheck}{HTML}{2E7D32}
\definecolor{labWarn}{HTML}{B23A48}

\lstset{
  backgroundcolor=\color{black!4},
  basicstyle=\ttfamily\footnotesize,
  breaklines=true,
  frame=leftline,
  framerule=2pt,
  rulecolor=\color{labAccent},
  numbers=left,
  numberstyle=\tiny\color{black!50},
  showstringspaces=false,
  tabsize=2
}

\setlength{\tabcolsep}{4pt}

% Let TeX add a little extra inter-word stretch as a last resort before
% it reports an Overfull \hbox. Only kicks in on lines that would
% otherwise overflow (long \texttt paths, URLs, CIDRs); never loosens a
% line that already fits. Keeps prose/itemize within the margin.
\setlength{\emergencystretch}{3em}

\pagestyle{fancy}
\fancyhf{}
\renewcommand{\headrulewidth}{0.4pt}
\lhead{\small \courseTitle{} -- Lab \labNumber}
\rhead{\small Maps to Section \labSection}
\cfoot{\thepage}

\newtcolorbox{stepbox}[1]{
  enhanced, breakable, arc=2pt, boxrule=0pt,
  colback=labAccent!7, colframe=labAccent, leftrule=3pt,
  fonttitle=\bfseries\small, coltitle=white,
  colbacktitle=labAccent,
  title={\faTools\ Step #1}
}

\newtcolorbox{warnbox}{
  enhanced, breakable, arc=2pt, boxrule=0pt,
  colback=labWarn!7, colframe=labWarn, leftrule=4pt,
  fonttitle=\bfseries\small, coltitle=white, colbacktitle=labWarn,
  title={\faExclamationTriangle\ Caution}
}

\newtcolorbox{notebox}{
  enhanced, breakable, arc=2pt, boxrule=0pt,
  colback=labAccent!4, colframe=labAccent, leftrule=2pt,
  fonttitle=\bfseries\small, coltitle=white, colbacktitle=labAccent,
  title={\faInfoCircle\ Note}
}

\newtcolorbox{goalbox}{
  enhanced, breakable, arc=2pt, boxrule=0pt,
  colback=labCheck!7, colframe=labCheck, leftrule=3pt,
  fonttitle=\bfseries\small, coltitle=white, colbacktitle=labCheck,
  title={\faBullseye\ Goal}
}

\newtcolorbox{deliverbox}{
  enhanced, breakable, arc=2pt, boxrule=0pt,
  colback=labSec!10, colframe=labSec, leftrule=3pt,
  fonttitle=\bfseries\small, coltitle=white, colbacktitle=labSec,
  title={\faClipboardList\ Deliverable}
}

% --- Solution-sheet boxes (used only by lab-solutions.tex) ----------
\newtcolorbox{solutionbox}[1]{
  enhanced, breakable, arc=2pt, boxrule=0pt,
  colback=labCheck!7, colframe=labCheck, leftrule=3pt,
  fonttitle=\bfseries\small, coltitle=white, colbacktitle=labCheck,
  title={\faCheckCircle\ #1}
}

\newtcolorbox{rubricbox}{
  enhanced, breakable, arc=2pt, boxrule=0pt,
  colback=labSec!7, colframe=labSec, leftrule=3pt,
  fonttitle=\bfseries\small, coltitle=white, colbacktitle=labSec,
  title={\faBalanceScale\ Grading rubric}
}

\titleformat{\section}{\large\bfseries\color{labAccent}}{\thesection}{0.5em}{}

\newcommand{\labheader}{%
  \noindent
  \begin{tikzpicture}
    \fill[labAccent] (0,0) rectangle (\textwidth, -1.6cm);
    \node[anchor=north west, text=white, font=\bfseries\large] at (0.6cm, -0.4cm) {Lab \labNumber{} -- \labTitle};
    \node[anchor=south west, text=white, font=\small] at (0.6cm, -1.3cm) {\courseTitle{} \quad $\bullet$ \quad Maps to Section \labSection{} \quad $\bullet$ \quad \labTime};
    \node[anchor=north east, fill=labSec, text=white, font=\bfseries, inner sep=4pt, rounded corners=2pt]
      at ([xshift=-0.4cm, yshift=-0.4cm]\textwidth,0) {LAB \labNumber};
  \end{tikzpicture}
  \vspace{4mm}
}


\begin{document}
\labheader

\begin{goalbox}
Lecturer / TA reference for Master Lab 06: a worked technique profile for
\textbf{ELECTRUM / Sandworm} with real ATT\&CK-for-ICS \texttt{Txxxx} IDs, a
valid STIX 2.1 snippet, a worked Diamond Model of Industroyer2, a coverage
table, model PIRs, common student mistakes, and a grading rubric. Accept the
same depth for a student who chose \textbf{CHERNOVITE} or \textbf{VOLTZITE}
instead --- the IDs differ, the tradecraft does not.
\end{goalbox}

\begin{solutionbox}{Step 1--2 -- model ICS technique profile (ELECTRUM / Sandworm)}
The profile is \emph{dual-matrix}: enterprise techniques for the IT intrusion,
ICS techniques for the OT impact. The pivot is typically \textbf{T0822 External
Remote Services} or \textbf{T0866 Exploitation of Remote Services} into the OT
DMZ. Accept a profile of 10--20 ICS techniques; the load-bearing ones for
Sandworm's grid attacks are below.

\begin{center}\footnotesize
\begin{tabular}{@{}p{2.0cm}p{4.0cm}p{8.2cm}@{}}
\toprule
\textbf{Txxxx} & \textbf{Technique} & \textbf{Evidence (incident / report)} \\
\midrule
T0883 & Internet-Accessible Device   & Exposed remote-access pathways into the OT DMZ (2015 Ukraine). \\
T0822 & External Remote Services     & Stolen VPN/RDP credentials to reach the SCADA network (2015). \\
T0859 & Valid Accounts               & Operator credentials harvested in the IT intrusion, reused in OT. \\
T0855 & Unauthorized Command Message & Industroyer issues breaker-open commands over IEC-101/104, OPC. \\
T0836 & Modify Parameter             & Manipulation of protection / setpoints. \\
T0814 & Denial of Service            & Industroyer's OPC/SIPROTEC DoS module against protective relays. \\
T0816 & Device Restart/Shutdown      & Forced relay/RTU restart to deepen the outage. \\
T0809 & Data Destruction             & Wiper component (KillDisk / CaddyWiper) erasing engineering hosts. \\
T0879 & Damage to Property           & Intent toward physical damage (relay disable $\to$ equipment risk). \\
T0826 & Loss of Availability         & The realised consequence: de-energised substations. \\
\bottomrule
\end{tabular}
\end{center}

\textbf{Enterprise sub-list (pre-pivot):} spearphishing (T1566), valid accounts
(T1078), remote services / RDP (T1021), credential dumping (T1003), data
destruction via wiper on IT (T1485). \textbf{Pivot sentence:} the adversary
crossed IT$\to$OT by abusing legitimate remote-access services into the SCADA
DMZ (T0822), not by a novel exploit.

{\footnotesize\emph{Sources:} Dragos / ESET, \emph{CRASHOVERRIDE / Industroyer}
2017; ESET, \emph{Industroyer2}, 2022; CISA \emph{AA22-103A}; Dragos ELECTRUM
threat-group page.}
\end{solutionbox}

\begin{solutionbox}{Step 2 -- ICS vs.\ enterprise: the grading trap}
The most common error is citing enterprise IDs in an ICS layer. Mark the layer
down if \texttt{techniqueID}s start with \texttt{T1} in an \texttt{ics-attack}
domain. The correct answer recognises that Sandworm/ELECTRUM ran an
\emph{enterprise} intrusion (T1566, T1078, T1021, T1003) and only the OT-impact
phase uses ICS IDs (T0855, T0836, T0814, T0816). Full marks require both lists
plus a named pivot technique (T0822 or T0866).
\end{solutionbox}

\begin{solutionbox}{Step 3 -- valid STIX 2.1 snippet}
A minimal accepted bundle. Note the ATT\&CK ID lives in an
\texttt{external\_reference} with \texttt{source\_name = mitre-ics-attack}, the
TLP marking is the canonical OASIS UUID (auto-supplied by \texttt{TLP\_AMBER}),
and relationships reference objects by \texttt{id}, not by name.

\begin{lstlisting}[]
{
  "type": "bundle",
  "id": "bundle--a1b2c3d4-0006-4f00-9aaa-electrum0001",
  "objects": [
    {
      "type": "intrusion-set", "spec_version": "2.1",
      "id": "intrusion-set--0006e1ec-...-electrum",
      "name": "ELECTRUM",
      "description": "OT-capable activity group (Dragos).",
      "object_marking_refs": ["marking-definition--f88d31f6-486f-44da-b317-01333bde0b82"]
    },
    {
      "type": "attack-pattern", "spec_version": "2.1",
      "id": "attack-pattern--0006a855-...-cmdmsg",
      "name": "Unauthorized Command Message",
      "external_references": [
        {"source_name": "mitre-ics-attack", "external_id": "T0855"}
      ]
    },
    {
      "type": "indicator", "spec_version": "2.1",
      "id": "indicator--0006idc1-...-mbwrite",
      "name": "Modbus write burst from non-EWS host",
      "pattern": "[network-traffic:dst_port = 502]",
      "pattern_type": "stix",
      "valid_from": "2026-05-21T00:00:00Z"
    },
    {
      "type": "relationship", "spec_version": "2.1",
      "id": "relationship--0006rel1-...-uses",
      "relationship_type": "uses",
      "source_ref": "intrusion-set--0006e1ec-...-electrum",
      "target_ref": "attack-pattern--0006a855-...-cmdmsg"
    }
  ]
}
\end{lstlisting}

\textbf{Validation:} \texttt{stix2\_validator campaign.json} must return no
errors. Reject bundles with: missing \texttt{spec\_version}, hand-typed
\texttt{created} strings that break the validator, relationships referencing
names instead of IDs, or enterprise \texttt{source\_name = mitre-attack} on an
ICS technique (should be \texttt{mitre-ics-attack}).
\end{solutionbox}

\begin{solutionbox}{Step 4 -- model detection-coverage table}
Honest gaps are worth more than overclaimed coverage. A good submission flags at
least two gaps.

\begin{center}\footnotesize
\begin{tabular}{@{}p{1.6cm}p{4.4cm}p{3.6cm}p{4.0cm}@{}}
\toprule
\textbf{Txxxx} & \textbf{Detection} & \textbf{Telemetry} & \textbf{Coverage} \\
\midrule
T0855 & Zeek notice on unsolicited Modbus write & \texttt{modbus.log} / \texttt{notice.log} & Full (lab proves it) \\
T0836 & Same notice, write-register func codes  & \texttt{modbus.log} & Partial (value-context missing) \\
T0846 & Conn-count heuristic (host fan-out)     & \texttt{conn.log} & Partial (noisy) \\
T0859 & ---                                     & no IAM/auth log in lab & \textbf{Gap} \\
T0814 & ---                                     & no relay/OPC telemetry & \textbf{Gap} \\
\bottomrule
\end{tabular}
\end{center}

Teaching point: the Bachelor lab can demonstrate the Modbus-impact techniques
convincingly but has no identity or protective-relay telemetry, so
account-abuse and DoS-on-relay techniques are honest blind spots --- exactly the
kind of gap a real OT SOC must acknowledge.
\end{solutionbox}

\begin{solutionbox}{Step 5 -- worked Diamond Model (Industroyer2, Apr 2022)}
\begin{itemize}[leftmargin=*, nosep]
  \item \textbf{Adversary:} Sandworm / ELECTRUM (GRU Unit 74455); tradecraft
        continuity with the 2016 Industroyer and CaddyWiper operators.
  \item \textbf{Capability:} Industroyer2 --- a single-protocol (IEC-104)
        variant with hard-coded target parameters; deployed with CaddyWiper,
        ORCSHRED, SOLOSHRED. Maps to T0855, T0836, T0816, T0809.
  \item \textbf{Infrastructure:} prior IT foothold in the energy provider;
        staging on internal hosts (Type 2 --- victim-resident, not external C2),
        scheduled execution timed to a maintenance window.
  \item \textbf{Victim:} a Ukrainian high-voltage electrical substation
        operator; targeted consequence = loss of availability (de-energise
        substations), with wipers to slow recovery.
  \item \textbf{Meta-features:} \emph{timestamp} April 2022; \emph{phase}
        ICS Cyber Kill Chain Stage 2 (Develop $\to$ Test $\to$ Deliver $\to$
        Execute ICS attack); \emph{result} attack detected and disrupted before
        full impact; \emph{confidence} HIGH (vendor + national-CERT corroboration,
        code continuity with 2016).
\end{itemize}
\textbf{Confidence justification:} two independent technical analyses (ESET +
CERT-UA) plus binary lineage to the 2016 tool. A single contradicting data point
that would lower it: credible evidence of a distinct operator group reusing the
leaked toolkit.
\end{solutionbox}

\begin{solutionbox}{Step 6 -- model consequence-driven PIRs (electric transmission)}
PIRs are questions anchored to a process consequence, with a source and a
trigger. Sample, ranked:

\begin{enumerate}[leftmargin=*, nosep]
  \item \textbf{PIR-1:} Is ELECTRUM developing/deploying IEC-104 or OPC
        capability against \emph{our} relay/RTU families? \emph{(Consequence:}
        loss of control, breaker operation. \emph{Source:} Dragos/CISA + vendor
        PSIRT. \emph{Trigger:} any advisory naming our controller/relay family.)
  \item \textbf{PIR-2:} Are stolen remote-access credentials for our OT DMZ for
        sale or in use? \emph{(Loss of view/control via T0822/T0859. Source:}
        ISAC + access-broker monitoring. \emph{Trigger:} our domain/VPN appears
        in a broker listing.)
  \item \textbf{PIR-3:} Is wiper tooling (CaddyWiper-class) being staged against
        sector peers? \emph{(Recovery denial after impact. Source:} E-ISAC +
        peer sharing. \emph{Trigger:} confirmed wiper at a peer utility.)
  \item \textbf{PIR-4:} Is there new public PoC for IEC-101/104 command
        injection? \emph{(T0855 lowers attacker cost. Source:} CVE feeds, conf
        talks. \emph{Trigger:} working PoC published.)
  \item \textbf{PIR-5:} Are protective-relay DoS techniques (T0814) being
        observed regionally? \emph{(Safety-trip / equipment damage. Source:}
        vendor PSIRT + national CERT. \emph{Trigger:} advisory or peer report.)
\end{enumerate}
\textbf{Grading trap:} a PIR that says ``monitor for malware'' or ``track the
threat landscape'' is not consequence-driven --- it is a generic IT outcome with
no process consequence and no trigger. Mark those down.
\end{solutionbox}

\begin{solutionbox}{Step 7 -- model detection (stretch)}
Accept a syntactically valid Sigma or Zeek rule that (a) carries the
\texttt{Txxxx} tag, (b) names a realistic false positive (commissioning /
batch HMI writes), and (c) has a tuning note (allowlist the EWS/HMI source IP,
or scope to off-hours). The lab body's Modbus-write-burst Sigma rule
(\texttt{T0855}/\texttt{T0836}) is the reference; a Zeek
\texttt{event modbus\_message} handler raising a notice when a non-allowlisted
\texttt{id.orig\_h} issues write func codes 5/6/15/16 is equally acceptable.
\end{solutionbox}

\begin{solutionbox}{Common student mistakes}
\begin{itemize}[leftmargin=*, nosep]
  \item \textbf{Enterprise IDs in an ICS layer} --- \texttt{T1059} where
        \texttt{T0807} (Command-Line Interface) belongs; the giveaway is
        \texttt{T1} prefixes in an \texttt{ics-attack} layer.
  \item \textbf{Invalid STIX} --- relationships pointing at names not IDs;
        hand-typed timestamps; wrong \texttt{source\_name}; bundle never run
        through the validator.
  \item \textbf{Indicators with no provenance} --- pasting unverified ``C2 IPs''
        from a blog; require defanged, source-cited, or plausibly-synthetic
        indicators.
  \item \textbf{Diamond with empty vertices} --- ``Adversary: a hacker'' instead
        of the named group + sub-cluster; missing confidence rating.
  \item \textbf{PIRs that are statements, not questions} --- or questions with no
        consequence and no trigger.
  \item \textbf{Switching groups mid-lab} --- profile for ELECTRUM, Diamond for
        CHERNOVITE; the artefacts must be internally consistent.
  \item \textbf{Overclaimed coverage} --- mapping every technique to ``Zeek''
        with no honest gaps; the lab telemetry cannot detect identity or relay
        techniques.
\end{itemize}
\end{solutionbox}

\begin{rubricbox}
Total \textbf{10 pts}; pass mark 5.
\begin{itemize}[leftmargin=*, nosep]
  \item \textbf{Technique mapping accuracy (3 pts):} 10--20 ICS techniques with
        correct \texttt{T0xxx} IDs and evidence sentences (2); enterprise/ICS
        split + named IT$\to$OT pivot (1). Auto-deduct 1 pt for any \texttt{T1}
        ID inside the ICS layer.
  \item \textbf{STIX validity (2 pts):} \texttt{campaign.json} passes the
        validator with zero errors (1); contains intrusion-set + $\geq$2
        attack-patterns with ATT\&CK \texttt{external\_id}s + indicators +
        relationships + TLP marking (1).
  \item \textbf{Detection mapping (2 pts):} coverage table with a row per
        technique (1); at least the Modbus-write techniques mapped to a real
        detection and $\geq$2 honest gaps named (1).
  \item \textbf{Analytic rigour (2 pts):} Diamond Model with all four vertices
        populated from cited sources + confidence rating (1); 5--7
        consequence-driven, source- and trigger-bearing PIRs (1).
  \item \textbf{Intel-product clarity (1 pt):} the 2-page assessment is readable
        in five minutes, TLP-marked, and the layer + bundle load without edits.
\end{itemize}
\textbf{Auto-fail Part B} if the STIX bundle was never validated. \textbf{Auto-deduct
2 pts} if any TLP-restricted source material was redistributed outside the cohort.
\end{rubricbox}

\end{document}
